\chapter*{ВВЕДЕНИЕ}
\addcontentsline{toc}{chapter}{ВВЕДЕНИЕ}

Обеспечение экологической безопасности водных ресурсов является одной из наиболее приоритетных задач на сегодняшний день. Доклад ООН указывает на то, что глобальный спрос на пресную воду растёт на чуть менее, чем 1\% в год и уже к 2030 году разница между доступной к употреблению воды и спросом на неё достигнет 40\%. Поскольку происходит загрязнение поверхностных вод, в частности рек и пресноводных озёр, и подземных ввод в следствие происшествия стихийных бедствий и неконтроллируемого антропогенного воздействия, в результате чего осуществляется существенное изменение состава и качества воды, используемой для питьевого и хозяйственно-бытового водоснабжения, требуется непрерывный мониторинг состояния гидросферы. В этом контексте одними из маркеров состояния природной воды выступают неорганические анионы \cite{unwaterWaterProsperityPeace2024}.

Современные требования к контролю качества воды в городских и промышленных условия, а также охране окружающей среды, диктуют постоянный поиск наиболее совершенных методов количественного анализа природных вод. Учитывая широкий спектр этих аналитических потребностей, использование методов, позволяющих одновременно определять множество компонентов, имеет особенное значение. Использование атомной-эмиссионной спектроскопии имеет долгую историю для определения многих лёгких элементов, но несмотря на это, наибольшоий прорыв в отношении определения неорганических ионов произошёл после внедрения жидкостной хроматографии в 1980-х годах в анализ качества воды, так как существенно упростило совместное количественное определенние ионов в воде. В настоящее время эти методы широко и регулярно используются в анализе качества воды. Также с момента первых экспериментальных применений капилярного электрофореза в 1980-х годах привлекает к себе всё больше внимания, являясь одним из единственных методов, сравнимых по эффективности с ВЭЖХ, что подверждается большим количеством выпускаемыйх статей и обзоров и общих тенденций в использовании аналитического капиллярного электрофореза \cite{pobozyApplicationCapillaryElectrophoresis2021b}.

Таким образом, была поставлена цель выполенения курсовой работы - проведения комплексного исследования методологии определения содержания анионов в природных водах

Для выполнения этой цели были поставлены следующие задачи:
\begin{enumerate}[label=\arabic*., leftmargin=\parindent, nosep]
	\item Провести аналитический обзор современных методов анализа природных вод;
	\item Составить методологию проведения эксперимента с необходимым оборудованием;
	\item Проанализировать результаты применения метода на реальном объекте.
\end{enumerate}
